\documentclass[12pt,a4paper]{article}
\usepackage[utf8]{inputenc}
\usepackage[T1]{fontenc}
\usepackage[english]{babel}
\usepackage{lmodern}
\usepackage{microtype}
\usepackage{amsmath, amssymb, amsfonts}
\usepackage[margin=2.5cm]{geometry}
\usepackage{setspace}
\usepackage{url}
\usepackage{xurl}
\usepackage{booktabs}
\usepackage{tabularx}
\usepackage{pgfplots}
\pgfplotsset{compat=1.18}
\usetikzlibrary{babel}

\onehalfspacing

\title{\textbf{P1: The Stationary-String Framework --- A Unified FENE-P Model for Galactic, Hadronic and Laboratory Phenomenology} \\[0.3cm]
\large (A Degrees-of-Freedom Ledger, a Fragility Map, and a Forensic Audit of Its Own Errors)}
\author{Fahri Meral}
\date{October 20, 2026}

\begin{document}

\maketitle

\begin{abstract}
This paper presents the FENE-P flux-string framework developed across
Theorems I--X and Grand Proof I--III as a \emph{methods} paper that lets the
framework score its own claims. Its locked invariants
($\Lambda_{\mathrm{Meral}}^{(0)} = 0.07226943267767906$,
$\mu_s^{\mathrm{Meral}} = 0.20440882367556615$,
$\mathcal{C}_{\mathrm{univ}} = 1.6260379365521303$,
$S_{\mathrm{loss}} = 19.75114251524387$,
$\mathcal{I}_{\mathrm{shield}} = 2\sqrt{2}$) follow from closed forms and
numerical integrals; the verification infrastructure comprises 38 automated
tests and seven independent cross-check scripts (94 independent checks in
total, none of which imports the solver). Our central contribution is a
\emph{degrees-of-freedom ledger} (Table~\ref{tab:ledger}): for each of 17
published comparisons it states the observable, the literature value, the
model value, the free parameters entering, whether those parameters were
fitted to that datum, and a fragility measure. The ledger shows that of the
17 comparisons only 2 are fits, 10 are transfers of the locked core, 2 are
genuine predictions and 2 are structural; grouping correlated rows leaves 8
independent observable groups against 9 free degrees of freedom, i.e.\ an
evidence ratio of 0.889. The fragility map further shows that the glueball
mean is unstable: a 1\% shift of $\mathcal{C}_{\mathrm{univ}}$ moves it by
27.8\%, so the 1.56\% agreement is \emph{not} robust. A forensic audit
record is finally given, listing eight error classes found and corrected
within this work (a wrong self-adjoint extension, transposed source
formulas, literature bands that turned out to be a \emph{ghost parameter},
an $H_0$ dependence, a tautological lower bound, and others). The framework
is therefore presented not as a proven theory but as a quantitatively
auditable and self-limiting unified parametrisation.
\end{abstract}

\section{Scope and the limits of the claims}
This paper deliberately separates three levels of claim:
\begin{enumerate}
    \item \textbf{Algebraic and numerical claims (exact):} the closed-form
    invariants, the self-adjoint extensions of the FENE-P string, the
    $p = 2.0000$ convergence order of the cell-centred finite-volume
    discretisation, and the analytic Temple--Kato enclosure. These are
    delivered together with the verification infrastructure.
    \item \textbf{Parametrisation claims (not derived):} the transport of
    locks such as $\mathcal{C}_{\mathrm{univ}}$ to galactic, Yang--Mills
    and laboratory observables. These involve ansatz choices, and the ledger
    (Table~\ref{tab:ledger}) states which row depends on which choice.
    \item \textbf{Out of scope:} the claim that the Yang--Mills mass gap is
    fully solved, or that dark matter is replaced. This paper makes neither
    claim.
\end{enumerate}

\section{The locked invariants}
The closed-form global-integrity invariant (a $\mathrm{GL}$ character
integral, Euler-derived) and the dimensionless coupling, saturation ratio
and shield invariant derived from it are
\begin{equation}
\label{eq:core}
\Lambda_{\mathrm{Meral}}^{(0)} = \frac{28}{79}\cdot\frac{1}{2}
\left[\frac{163}{578} + \frac{64\sqrt{17}}{4913}
\ln\!\left(\frac{21+5\sqrt{17}}{4}\right)\right] = 0.07226943267767906,
\end{equation}
\begin{equation}
\label{eq:core2}
\begin{split}
\mu_s^{\mathrm{Meral}} &= 2\sqrt{2}\,\Lambda_{\mathrm{Meral}}^{(0)}
= 0.20440882367556615, \\
\mathcal{C}_{\mathrm{univ}} &= \frac{F_{\mathrm{base}}}{F_{\mathrm{eff}}}
= 1.6260379365521303, \\
S_{\mathrm{loss}} &= \frac{e^{2\sqrt{2}}}{1-\eta_{\mathcal{T}_0}}
= 19.75114251524387.
\end{split}
\end{equation}
\begin{equation}
\label{eq:core3}
\mathcal{I}_{\mathrm{shield}} = \ln\!\left[S_{\mathrm{loss}}(1-\eta_{\mathcal{T}_0})\right] = 2\sqrt{2}.
\end{equation}
These numbers carry the framework's ``no new calibration constant'' claim:
no external parameter enters the core.

\section{Verification infrastructure}
The solver \texttt{solver\_engine.py} contains 38 automated tests
(T1--T37). In addition, seven independent scripts re-derive the same
physical constants without ever importing the solver:
\texttt{mini\_check\_viii} (8/8), \texttt{mini\_check\_ix} (15/15),
\texttt{mini\_check\_x} (15/15), \texttt{mini\_check\_gp1} (15/15),
\texttt{mini\_check\_gp2} (17/17), \texttt{mini\_check\_gp3} (15/15) and the
series-wide \texttt{mini\_check\_audit} (9/9) --- 94 independent checks in
total. The audit script re-derives $\mathcal{C}_{\mathrm{univ}}$ with an
ODE integration instead of the Gauss--Legendre quadrature and finds
agreement to $5.6\times10^{-13}$; the identity
$\Pi_{\mathrm{Meral}} = \mu_s^{\mathrm{Meral}}$ is confirmed to
$5.6\times10^{-17}$.

\section{The degrees-of-freedom ledger}
Table~\ref{tab:ledger} lists every published comparison against four
questions. The \emph{evidence class} means: \textsc{Fit} (a parameter was
chosen to reproduce this datum), \textsc{Fit-like} (a multiplicative factor
was added after a deviation was seen), \textsc{Transfer} (the locked core is
carried over from another domain), \textsc{Genuine} (a discriminating
prediction with no free parameter) and \textsc{Structural} (a rule rather
than a number).

\begin{table}[!htbp]
\centering\scriptsize
\setlength{\tabcolsep}{2.5pt}
\renewcommand{\arraystretch}{1.05}
\caption{\textbf{Degrees-of-freedom ledger: evidence class and deviation for
17 published comparisons.} Rows are cross-checked against the JSON emitted
by \texttt{ledger.py} (tolerance $2\times10^{-3}$).}
\label{tab:ledger}
\vspace{0.1cm}
\begin{tabularx}{\textwidth}{@{}lXrrll@{}}
\toprule
Id & Observable & Literature & Model & Dev. & Class \\
\midrule
T15a & a1S/h (1.03 T) [kHz] & $1.4204048\times 10^{6}\pm 5.6$ & $1.4204048\times 10^{6}$ & 0.00\% & Fit \\
T15b & a-b-g (1.03 T) [g] & $0.75\pm 0.16$ & $0.75$ & 0.00\% & Fit \\
T19 & read-out frequency (993 mT) [kHz] & $453$ & $453$ & 0.00\% & Transfer \\
T22 & Ps band budget ratio $R$ & $3.3387$ & $3.3387$ & 0.00\% & Genuine \\
T25 & loss suppression factor & $31.4$ & $31.41$ & 0.03\% & Transfer \\
T29 & $a_{\mathrm{Meral}}$ [m/s$^2$] & $1.20\times 10^{-10}\pm 2.8\times 10^{-12}$ & $1.1466\times 10^{-10}$ & 4.45\% & Transfer \\
T29b & BTFR log-log slope & $4$ & $4$ & --- & Structural \\
T30 & 5-galaxy mean deviation [\%] & target 0 & $4.5375$ & 4.54\% & Transfer \\
T31 & $d_{\mathrm{ring}}/\theta_g$ & $10\pm 0.4$ & $10.002$ & 0.02\% & Transfer \\
T32 & $\Delta m_{\mathrm{gap}}/\sqrt{\sigma}$ & $3.405\pm 0.021$ & $3.3914$ & 0.40\% & Fit-like \\
T33 & 4-glueball mean deviation [\%] & target 0 & $1.5563$ & 1.56\% & Fit-like \\
T34a & Cornell coefficient & $0.295\pm 0.010$ & $0.2985$ & 1.18\% & Transfer \\
T34b & $c\bar{c}$ $2S-1S$ [MeV] & $589.19$ & $596.89$ & 1.31\% & Transfer \\
T34c & $b\bar{b}$ $2S-1S$ [MeV] & $562.96$ & $522.25$ & 7.23\% & Transfer \\
T36 & $\Delta\nu/\nu$ (2004 event) & $2.08\times 10^{-6}$ & $2.0711\times 10^{-6}$ & 0.43\% & Transfer \\
T35 & resistance envelope lower edge & $12$ & $12.147$ & 1.22\% & Transfer \\
T37 & $O_{\mathrm{exact}}$ band ratio & $0.08414$ & $0.084136$ & 0.01\% & Genuine \\
\bottomrule
\end{tabularx}
\end{table}

\subsection{Three results the ledger makes explicit}
\paragraph{(i) The evidence ratio.} Grouping rows that share a knob set
leaves 8 \emph{independent} observable groups, against 7 ansatz plus 2 fit
parameters, i.e.\ 9 free degrees of freedom. The evidence ratio is
therefore $8/9 = 0.889$. The framework cannot carry a predictive burden in
its present form: the agreement is a natural consequence of the model's
flexibility, and the ledger says so explicitly.
\paragraph{(ii) The fragility map.} The most sensitive row is T33: a 1\% shift
of the locked $\mathcal{C}_{\mathrm{univ}}$ moves the mean glueball
deviation by 27.8\%, because the $0^{*++}$ factor multiplies
$\mathcal{C}_{\mathrm{univ}}$ while the gap scale does as well. The 1.56\%
agreement is therefore \emph{not} robust. By contrast T34b (charmonium,
0.05\%) and T36 (Vela, 0.17\%) are practically stable, because
$\mu_s^{\mathrm{Meral}}$ is fixed to 14 digits.
\paragraph{(iii) Only two genuine predictions.} Of 17 comparisons, only two
(T22 and T37) contain no free parameter and are genuinely discriminating;
both live on the laboratory side. The remaining fifteen are fits or
transfers. The same table therefore exposes the strongest and the weakest
parts of the series at once.

\section{Forensic audit: error classes found and fixed}
The second contribution of this work is the record of the framework's own
errors.
\begin{enumerate}
    \item \textbf{Wrong self-adjoint extension (T32).} The
    $u=\sqrt{x}\psi$ transformation selected an incorrect extension at the
    critical $-1/(4x^2)$ coupling, violating the variational upper bound
    ($5.9158 > 5.6468$). Replacing it with a cell-centred finite-volume
    scheme gives $p = 1.99995$ and
    $\lambda_0^{\mathrm{cont}} = 5.6102039939$; an independent shooting
    integration gives $5.6102039940$.
    \item \textbf{Tautological lower bound (T32).} The Temple--Kato lower
    bound had been derived from the numerical matrix' own eigenvector. It
    was replaced by an analytic, grid-independent enclosure,
    $\lambda_0 \in [5.4807, 5.7674]$.
    \item \textbf{Transposed source formulas (T37).} The published form of
    \texttt{arXiv:2605.21289} Eq.~(12) and the substituted variant had
    been interchanged; corrected to $159.37$ and $9.31$ respectively.
    \item \textbf{A literature band that was a ghost parameter (T35).}
    The bands $r_{\mathrm{ve}}/r_{\mathrm{N}} \in [1.55,1.70]$ and
    $\Delta h_{\mathrm{ve}}/\Delta h_{\mathrm{N}} \in [0.30,0.35]$ entered
    this work without literature provenance --- git archaeology shows they
    arrived with the present series --- and do not occur in the cited
    source; the claim was retracted.
    \item \textbf{$H_0$ dependence (T29).} $a_{\mathrm{Meral}}$ is directly
    proportional to $H_0$; the comparison was reframed as a test of the
    dimensionless ratio.
    \item \textbf{Post-hoc multiplicative factors (T32/T33).} The chiral
    dressing $(1+\Lambda_{\mathrm{Meral}}^{(0)}/3)$ and the topological
    $0^{-+}$ term were added after deviations were observed; both are
    labelled as extra model interpretations.
    \item \textbf{A model contradicting its own source (T35).} The
    single-mode laminar ansatz could not reproduce the measured flow
    resistance; a two-regime barrier bridge was constructed, and the fact
    that this construction was made \emph{after} the measurement is stated
    explicitly.
    \item \textbf{Hard-coded ``independent'' checks.} Four mini scripts
    copied $\mathcal{C}_{\mathrm{univ}}$ as a constant; an independent
    re-derivation (\texttt{mini\_check\_audit}) was added.
\end{enumerate}

\section{What would upgrade the claims}
The predictive value of the framework grows only if the ansatz choices are
derived. Three concrete tasks follow: (i) derive the effective potential
$V_{\mathrm{eff}}$ from the flux-string dynamics instead of postulating
it; (ii) derive the resistance/core-broadening link through the locked
$S_{\mathrm{loss}}$; (iii) produce at least three quantitative predictions
in domains never used to fit (for example kaon spectra, or transient
superfluid cores). The third is the kind of output the field rewards.

\section{Reproducibility}
The code lives with Theorems I--X in the \texttt{stationary-strings-v}
repository (under \path{duran_sicimler_sim/}). Running \texttt{python}
\path{solver_engine.py} executes the full test suite; \texttt{python}
\path{ledger.py} regenerates the degrees-of-freedom ledger, cross-checks it
against the solver output at $10^{-9}$ tolerance and verifies the table
in this paper against the regenerated values; \path{mini_check_audit.py}
re-derives the locked chain independently. All packages are released under
CC BY 4.0.

\begin{thebibliography}{99}
    \bibitem{meral5} Meral, F. (2026c). \textit{Stationary Strings
    Theorem V}. Open-source solver package: \texttt{stationary-strings-v}
    (\url{https://github.com/fahrim/stationary-strings-v}).
    \bibitem{meral12} Meral, F. (2026). \textit{Grand Proof I--III}.
    Open packages: \texttt{grand-proof-i}, \texttt{grand-proof-ii},
    \texttt{grand-proof-iii}.
    \bibitem{al26} Akbari, R. \textit{et al.} (2026). Four ppm measurement of
    the antihydrogen ground-state hyperfine splitting. \textit{Nature}
    \textbf{653}, 1022--1026.
    \bibitem{base26} Leonhardt, M., Morgner, J. \textit{et al.} (2026). Road
    transport of trapped antiprotons. \textit{Nature} \textbf{657}, 906--911.
    \bibitem{mcg16} McGaugh, S.~S., Lelli, F. \& Schombert, J.~M. (2016).
    Radial acceleration relation in rotationally supported galaxies.
    \textit{Phys.\ Rev.\ Lett.\ \textbf{117}}, 201101.
    \bibitem{gs00} Groisman, A. \& Steinberg, V. (2000). Elastic turbulence in
    a polymer solution flow. \textit{Nature} \textbf{405}, 53--55.
    \bibitem{link99} Link, B., Epstein, R.~I. \& Lattimer, J.~M. (1999).
    Pulsar constraints on neutron star structure.
    \textit{Phys.\ Rev.\ Lett.\ \textbf{83}}, 3362--3365.
    \bibitem{sy26} Siegert, T. \& Yoneda, H. (2026). On the possibility of an
    extragalactic positron annihilation signal. arXiv:2605.21289.
\end{thebibliography}

\end{document}
